%! TeX root: ../article.tex
\section{Metodología}
\label{sec:metodologia}

\subsection{Conjunto de datos}
\label{subsec:dataset}
El banco de imágenes de grafiti proporcionado para este trabajo se compone de
$6681$ fotografías capturadas en condiciones reales en la ciudad de Salamanca,
tomadas entre $2022$ y $2023$. De manera adicional, se dispone de un conjunto
de datos adicional de $1106$ imágenes de la ciudad de Cuenca, cedidas por la
iniciativa StopGrafiti\footnote{\url{https://stopgrafiti.com/}}, que se emplea
para el ajuste fino de algunos componentes del sistema.

Ninguno de los dos bancos de datos cuenta con anotaciones de ningún tipo, lo
que impide la evaluación supervisada directa de los modelos. Además, la
existencia de varios grafitis por imagen hace que la anotación a nivel de
imagen no sea suficiente para evaluar la calidad de los clústeres. Por tanto,
se opta por recortar y etiquetar manualmente varios subconjuntos de grafitis,
de modo que la evaluación supervisada se realice sobre estos recortes
etiquetados.

Se etiquetan manualmente recortes por estilo, asignándoles una de las cuatro
categorías estilísticas (firmas rápidas, \textit{tags}; letras redondeadas con
no más de dos colores, \textit{throw-ups}; texto mucho más elaborado,
\textit{pieces}; y personajes animados u otros elementos que no sean texto,
\textit{characters}), y otros recortes por autoría. La procedencia de cada
subconjunto difiere: el entrenamiento de estilo contiene recortes del banco
principal, mientras que el de autoría se restringe a StopGrafiti; el reparto y
su justificación se detallan al describir la evaluación extrínseca
(\ref{subsubsec:eval-extrinseca}). El reparto por clase de los recortes
etiquetados por estilo es desigual: las categorías \textit{tag}, \textit{piece}
y \textit{throw-up} concentran la mayor parte de los recortes, mientras que
\textit{character} resulta claramente minoritaria, lo que condiciona la
interpretación posterior de las métricas extrínsecas por estilo. Por último, se
utilizan imágenes de ambos grupos para el ajuste fino del detector de grafitis.

De las $6681$ imágenes de Salamanca, $265$ se reservan para entrenamiento
(ajuste fino del detector y de las cabezas, así como los recortes de autoría),
y las $6416$ restantes constituyen el conjunto de evaluación sin supervisión
del sistema completo, del que nunca se extraen recortes de entrenamiento. La
tabla~\ref{tab:subconjuntos} resume los subconjuntos etiquetados empleados, su
propósito, procedencia y tamaño.

\begin{table}[h!]
  \centering
  \begin{tabular}{|l|l|l|r|} \hline
    \textbf{Subconjunto} & \textbf{Propósito} & \textbf{Procedencia} & \textbf{Tamaño} \\ \hline
    Estilo (entr.) & Cabezas de estilo & StopGrafiti + Salamanca & $385$ \\ \hline
    Estilo (eval.) & M. Extrínsecas & Banco completo & $294$ \\ \hline
    Autoría (entr.) & Cabezas de autoría & StopGrafiti & $321$ \\ \hline
    Autoría (eval.) & M. Extrínsecas & Salamanca & $185$ \\ \hline
    Detección & Ajuste del detector & Ambos bancos & $447$ \\ \hline
    No superv. & Agrupamiento & Salamanca (disjunto) & $6416$ \\ \hline
  \end{tabular}
  \caption[Subconjuntos etiquetados empleados]{Subconjuntos empleados (en recortes, salvo detección y no
  supervisión, en imágenes). Los recortes de entrenamiento nunca se solapan con
  los de evaluación.}
  \label{tab:subconjuntos}
\end{table}

\subsection{Visión general del flujo de trabajo}
\label{subsec:vision}
El sistema implementado procesa cada fotografía en cuatro etapas operativas:
segmentación de los grafitis contenidos en la imagen, extracción de un vector
de características por recorte, persistencia del vector junto a sus metadatos
en un almacenamiento estructurado y, en un segundo tiempo, recuperación masiva
del conjunto almacenado para reducirlo y agruparlo (figura~\ref{fig:pipeline}).

\begin{figure}[h!]
  \centering
  \shorthandoff{<>}%
  \begin{tikzpicture}[
    >=Stealth, node distance=11mm and 17mm,
    block/.style={draw, rounded corners, minimum height=9mm, minimum width=22mm,
                  inner sep=5pt, align=center, font=\small},
    store/.style={block, fill=green!12},
    io/.style={font=\footnotesize, align=center},
    lbl/.style={font=\scriptsize, midway, fill=white, inner sep=2pt},
  ]
    % ingesta
    \node[io] (img) {fotografía};
    \node[block, right=of img] (seg) {Segmentador};
    \node[block, right=of seg] (ext) {Extractor};
    \node[store, right=of ext] (sto) {Almacenamiento};

    \draw[->] (img)--(seg);
    \draw[->] (seg)--(ext) node[lbl, above]{recortes};
    \draw[->] (ext)--(sto) node[lbl, above]{\textit{embeddings}};

    % agrupamiento
    \node[block, below=14mm of seg] (red) {Reducción};
    \node[block, right=of red] (clu) {Agrupamiento};
    \node[io, right=of clu] (out) {clústeres\\y métricas};

    \draw[->] (red)--(clu);
    \draw[->] (clu)--(out);

    % leer almacenamiento
    \draw[->] (sto.south) -- ++(0,-7mm) -| (red.north)
        node[lbl, pos=0.25, above]{carga masiva};

    \node[draw, dotted, fit=(img)(seg)(ext), inner sep=3mm,
          label=above:{\scriptsize Fase de ingesta (incremental)}] {};
    \node[draw, dotted, fit=(red)(clu)(out), inner sep=3mm,
          label=below:{\scriptsize Fase de agrupamiento (puntual)}] {};
  \end{tikzpicture}%
  \shorthandon{<>}%
  \caption[Flujo de trabajo en dos fases]{Flujo de trabajo en dos fases. El almacenamiento es el punto de
  acoplamiento entre la ingesta incremental y la fase de agrupamiento puntual.}
  \label{fig:pipeline}
\end{figure}


Cada una de las cinco etapas es un componente intercambiable tras una interfaz
común; el catálogo completo de opciones disponibles para cada una (modelos de
extracción, algoritmos de agrupamiento, técnicas de reducción y
\textit{backends} de almacenamiento) se recoge en el
anexo~\ref{subsec:configuraciones}. Este desacoplamiento es el que habilita la
evaluación combinatoria que forma el trabajo: cada combinación define una
configuración independiente y reproducible cuyo coste y calidad pueden
compararse con cualquier otra.

\subsection{\textit{Pipeline} operativo}
\label{subsec:pipeline}

\subsubsection{Ingesta y generación de \textit{embeddings}}
La extracción de características de cada fotorgraía comienza con la ingesta de
esta imagen por el segmentador escogido. Que devuelve una lista con las cajas
delimitadoras de los grafitis presentes en la imagen. Estas cajas se fusionan
según criterios de solapamiento, proximidad y confianza, y se recortan sobre la
imagen original para obtener los grafitis individuales. Cada recorte se entrega
al modelo de extracción de características, que produce un vector de
características de dimensión fija que representa el grafiti. Finalmente, el
vector resultante, junto con la referencia a la fotografía original y las
coordenadas del recorte se almacenan en el sistema elegido para su posterior
recuperación masiva.

En el caso de que no se utilice segmentador se devuelve una única caja que
abarca toda la imagen, de modo que el sistema se comporta como un agrupador de
imágenes completas. Este comportamiento también se puede habilitar o
deshabilitar para los modelos de detección de caso que no detecten ningún
grafiti.

\subsubsection{Recuperación, reducción y agrupamiento}
Una vez completada la ingesta, la fase de agrupamiento opera sobre el contenido
completo del almacenamiento. En primer lugar, se recuperan todos los
\textit{embeddings} junto con sus metadatos en una única consulta. A
continuación, si la configuración declara una reducción, se aplica sobre cada
vector de características el método de reducción seleccionado, obteniendo un
nuevo conjunto de vectores. Finalmente, el algoritmo de agrupamiento se ajusta
sobre el conjunto reducido, produciendo una etiqueta de clúster por cada
vector. A partir de esta información se calculan las métricas de calidad
intrínsecas y extrínsecas, y se conserva el mapeo entre identificadores de
clúster y listas de recortes para la generación de visualizaciones
interactivas. Esta secuencia de pasos se ejecuta además varias veces para cada
configuración para obtener estadísticas sobre la variabilidad de las métricas
de coste y calidad.

\subsection{Procedimiento experimental}
\label{subsec:experimental}
La evaluación se articula en dos ejes. El primero, y principal, es el
rendimiento computacional, cuantificado mediante los tiempos absolutos de cada
etapa del \textit{pipeline} (preparación, ingesta, reducción, agrupamiento y
búsqueda por similitud), el \textit{throughput} de ingesta en imágenes por
segundo, la latencia de búsqueda por consulta y la escalabilidad de estas
magnitudes con el tamaño del banco; todas ellas son métricas objetivas e
independientes de la existencia de etiquetas. El segundo eje, secundario, es la
calidad de los clústeres resultantes, evaluada con las métricas internas y
extrínsecas definidas en \ref{subsubsec:eval-intrinseca} y
\ref{subsubsec:eval-extrinseca}.

\subsubsection{Definición del espacio de configuraciones}
\label{subsubsec:espacio-configs}
Cada experimento queda descrito por un fichero de configuración en formato
JSON. El fichero contiene una lista de experimentos. Cada experimento declara
un identificador único y cinco componentes con la forma \texttt{\{type,
params\}}: \texttt{storage}, \texttt{embedding}, \texttt{segmenter},
\texttt{reduction} y \texttt{clustering}, además de un bloque opcional
\texttt{similarity\_search}. Esta forma declarativa garantiza que toda la
información necesaria para reproducir un experimento queda en un único fichero
versionable.

Para barrer sistemáticamente el espacio combinatorio se utiliza un fichero
auxiliar en el que cada componente se sustituye por una lista de alternativas.
Un \textit{script} expande este fichero mediante producto cartesiano,
materializando un experimento en el archivo final por cada combinación válida.

\subsubsection{Protocolo de escalabilidad}
\label{subsubsec:escalabilidad}
La caracterización del coste computacional en función del tamaño del banco
constituye uno de los ejes centrales del trabajo. Para obtenerla, se ejecuta
cada configuración seleccionada sobre una sucesión creciente de tamaños $N$
controlada mediante el parámetro \texttt{--limit}. La selección de tamaños es
reducida en la mayoría de experimentos para evitar que el coste total de la
evaluación se dispare, pero se realiza un barrido completo sobre un rango de
tamaños lo suficientemente amplio como para identificar tendencias de
crecimiento.

El muestreo de las imágenes para cada $N$ es determinista: las fotografías se
ordenan por su nombre de fichero y se toman las primeras $N$. Para cada tamaño
se obtienen las métricas temporales descritas en \ref{subsubsec:cronometria},
lo que permite caracterizar el crecimiento de cada etapa por separado y, a
partir de ahí, estimar empíricamente el órden de crecimiento de la función de
coste respecto a $N$ para cada etapa.

\subsubsection{Reutilización de almacenamientos}
\label{subsubsec:reutilizacion}
La ingesta es, en la mayoría de ocasiones, la etapa más costosa del
\textit{pipeline}. Para evitar repetirla, se crea un identificador para cada
almacenamiento a partir de la configuración de ingesta (segmentador, modelo de
representación y parámetros asociados) que luego es usado para situar la base
de datos en un directorio específico (o nombrar la tabla en el caso de
PostgreSQL). Cuando el sistema detecta que el almacenamiento ya existe y no se
ha declarado explícitamente que se limpie (\texttt{clear\_storage: true}), la
fase de ingesta se omite y se reutiliza el almacenamiento existente. Las
métricas de ingesta se almacenan en la tabla de metadatos del almacenamiento
para que, en casos como estos, se reporten los tiempos reales de ingesta en
lugar de los tiempos de caché, que no reflejarían el coste real de la etapa.

\subsubsection{Instrumentación y cronometría}
\label{subsubsec:cronometria}

\begin{sloppypar}
    El \textit{runner} de evaluación mide cuatro etapas: ingesta, reducción,
    agrupamiento y búsqueda por similitud. Cada una se envuelve en un gestor de
    contexto (\texttt{profile\_stage}) que registra los instantes de entrada y
    salida con \texttt{time.perf\_counter}, un reloj monotónico de alta resolución.
    La diferencia entre ambos instantes es el tiempo de pared (\textit{wall time})
    de la etapa, en segundos y con resolución de microsegundos. La preparación de
    componentes y la limpieza del almacenamiento quedan fuera del gestor de
    contexto y no se contabilizan.
\end{sloppypar}

Para acotar la variabilidad entre ejecuciones, las etapas de reducción y
agrupamiento se repiten $K$ veces dentro del mismo proceso, con las semillas
$0, 1, \ldots, K-1$. La semilla afecta a la reducción estocástica (UMAP) y a
los agrupamientos estocásticos (K-medias, GMM, \textit{spectral clustering} y
\textit{affinity propagation}); los algoritmos deterministas (DBSCAN, OPTICS,
HDBSCAN y aglomerativo) la reciben por la interfaz común pero la ignoran. Cuando
la reducción y el agrupamiento son ambos deterministas, las repeticiones serían
idénticas, así que se fuerza $K=1$; esas configuraciones, en consecuencia, no
aportan barras de error.

Cuando $K \geq 2$, la primera repetición ($i=0$) se descarta porque absorbe los
efectos de inicialización (compilación \textit{just-in-time} de operadores y
poblamiento de cachés del sistema operativo); sobre las $K-1$ restantes se
reportan la media aritmética y la desviación típica de cada métrica de coste.
La calidad del agrupamiento, en cambio, se evalúa sobre las $K$ repeticiones
completas: no la afectan los costes de inicialización y así se cuantifica la
sensibilidad de las métricas intrínsecas al muestreo aleatorio.

La ingesta y la búsqueda por similitud se miden una sola vez por ejecución. La
ingesta es determinista y, además, la reutilización de almacenamiento
(\ref{subsubsec:reutilizacion}) exige una cronometría única; la búsqueda por
similitud es barata y su variabilidad entre configuraciones se observa
directamente a lo largo del barrido de tamaños. Los experimentos se ejecutan
sobre un sistema sin otras cargas significativas (\ref{subsec:entorno}) para
minimizar interferencias externas.

\subsubsection{Selección automática de hiperparámetros}
\label{subsubsec:hiperparam-auto}
Salvo que se desee evaluar explícitamente el efecto de fijar un hiperparámetro
a un valor concreto, todas las configuraciones del \textit{benchmark} emplean
por defecto la estimación automática del hiperparámetro principal de K-medias,
DBSCAN y los modelos de mezcla gaussiana, mediante el método del codo, la curva
de $k$-distancias y el criterio bayesiano de información respectivamente. En los
tres casos las semillas aleatorias se fijan para que el resultado sea
reproducible.

La elección de estos algoritmos también atiende a su coste computacional:
K-medias cuenta con variantes escalables como la versión
\textit{mini-batch}~\cite{sculley_web-scale_2010}, mientras que los métodos
basados en densidad presentan un crecimiento superlineal con el tamaño del
banco. En cualquier caso, el rendimiento en calidad depende fuertemente del
conjunto de datos~\cite{franti_k-means_2018,ahmed_k-means_2020}, lo que
refuerza el enfoque comparativo adoptado en este trabajo.

Cuando una configuración no fija \texttt{n\_clusters} y emplea varias
repeticiones, el número de clústeres autodetectado en la primera repetición se
propaga al resto, de modo que las repeticiones varían únicamente por la semilla
y las métricas agregadas describen una sola familia de particiones en lugar de
mezclar particiones con distinto número de clústeres. El número de clústeres
obtenido en cada repetición se registra para poder verificar que la propagación
se mantiene.

\subsubsection{Protocolo de búsqueda por similitud}
\label{subsubsec:similarity-search}
La medición del coste de búsqueda por similitud se realiza sobre los
\textit{backends} de almacenamiento descritos anteriormente. Tras realizarse la
ingesta, se seleccionan $n$ \textit{embeddings} de manera aleatoria (pero
reproducible) y se buscan los $k$ vecinos más próximos a cada uno de ellos,
midiendo el tiempo total de las consultas. La distancia utilizada para la
búsqueda puede ser tanto euclidiana como coseno, y el número de vecinos también
es configurable.

El tiempo total de las consultas se mide con el mismo procedimiento que el
resto de etapas. De la misma forma que en el almacenamiento, la búsqueda por
similitud es determinista una vez fijada la semilla que determina las entradas
que buscar, así que no se repite varias veces como la reducción y el
agrupamiento. El coste de esta etapa se espera que crezca como $O(n^2)$ para el
recorrido lineal de \texttt{SQLiteStorage} y como $O(n \log n)$ para el índice
HNSW de \texttt{PostgreSQLStorage}.

De esta etapa también se reporta la distancia media a los $k$ vecinos
recuperados, promediada sobre todas las consultas. Si el \textit{backend}
dispone de índice ANN (como es el caso de \texttt{PostgreSQL} con
\texttt{HNSW}), también se reporta el R@$k$ de sus vecinos frente
a una búsqueda exacta por fuerza bruta, que cuantifica la pérdida de exactitud
del índice aproximado y es nulo en los \textit{backends} exactos.

\subsubsection{Evaluación intrínseca de los agrupamientos}
\label{subsubsec:eval-intrinseca}
La calidad intrínseca se cuantifica con tres métricas que miden la cohesión
intracluster y la separación intercluster directamente sobre los
\textit{embeddings}: el coeficiente de silueta~\cite{rousseeuw_silhouettes_1987},
el índice de Calinski--Harabász~\cite{calinski_dendrite_1974} y el índice de
Davies--Bouldin~\cite{davies_cluster_1979}. Se calculan de
manera operativa sobre los \textit{embeddings} originales en su espacio de
partida, no sobre los \textit{embeddings} reducidos utilizados por el
agrupamiento. Esto se hace porque el objetivo de estas métricas es evaluar la
calidad del agrupamiento de los datos originales, no la habilidad del reductor
para deformar el espacio de modo que los clústeres aparezcan más compactos o
separados.

Los algoritmos basados en densidad (DBSCAN, OPTICS y HDBSCAN) pueden asignar la
etiqueta $-1$ a puntos considerados ruido. Estos puntos se excluyen del cálculo
del coeficiente de silueta y de los índices de Calinski--Harabász y
Davies--Bouldin, lo que evita penalizar a los algoritmos que identifican y
reportan ruido. Por otro lado, un agrupamiento que rechace una gran fracción de
puntos como ruido podría quedar artificialmente premiado por agrupar solo los
puntos muy cercanos entre sí; por ello estas métricas se complementan con la
proporción de puntos etiquetados como ruido y la variedad de tamaños de los
clústeres resultantes.

\subsubsection{Evaluación extrínseca sobre los subconjuntos etiquetados}
\label{subsubsec:eval-extrinseca}
La evaluación supervisada hace uso de métricas extrínsecas (ARI, NMI y F1 por
pares, definidas más adelante) y opera sobre dos conjuntos de recortes
etiquetados, uno por cada tarea supervisada del trabajo: evaluación por estilo
y evaluación por autor. En ambos casos, se utilizan los recortes directamente
como imágenes a evaluar en vez de aportar la imagen original y utilizar un
segmentador para extraer los recortes de evaluación, porque el objetivo es
medir la calidad del espacio de representación para agrupar grafitis, no la
capacidad de un sistema completo. Tampoco se utilizan imágenes completas,
puesto que muchas de estas contienen varios grafitis de estilos y autores
distintos, haciendo difícil atribuir una etiqueta de estilo o autor a la imagen
completa y complicando la interpretación de las métricas.

Ambas evaluaciones comparten métricas y procedimiento, pero responden a
preguntas distintas, y esa diferencia justifica el reparto asimétrico de las
anotaciones entre entrenamiento y evaluación. Las etiquetas de estilo definen
una tarea de clasificación cerrada sobre cuatro categorías fijas, así que basta
con que no se solapen las fotografías entre entrenamiento y evaluación. La
diversidad estilística se cubre mejor agregando recortes de Salamanca y de
Cuenca (el banco StopGrafiti) al entrenamiento, y el conjunto de evaluación,
extraído del banco principal de Salamanca, es disjunto a nivel de fotografía.
La evaluación por estilo mide así el grado en que el espacio de representación
agrupa grafitis estilísticamente homogéneos.

Las etiquetas de autor, en cambio, definen un problema abierto de aprendizaje
de métricas: en el uso real del sistema, los autores de los grafitis a agrupar
son en su mayoría desconocidos para el modelo en el momento de la inferencia.
Si las cabezas de autoría se entrenasen sobre los mismos identificadores de
autor que luego aparecen en la evaluación (aun utilizando fotografías
distintas), las métricas medirían la capacidad de reidentificar autores ya
vistos, no la de construir un espacio en el que autores nuevos queden
separados. Para evitar este sesgo se sigue la convención habitual de la
verificación de identidad (FaceNet~\cite{schroff_facenet_2015} y la literatura
de reidentificación de personas): el conjunto de evaluación contiene únicamente
identificadores de autor disjuntos de los vistos durante el entrenamiento.

Esta restricción se cumple entrenando las cabezas de autoría exclusivamente con
los recortes de Cuenca y evaluándolas sobre los de Salamanca. Dado que ambas
ciudades están geográficamente separadas, es razonable asumir que los conjuntos
de autores no se solapan. Esto, aunque no se ha verificado formalmente, se ha
comprobado de manera informal al realizar la segmentación y anotación de los
recortes, y no se ha encontrado ningún caso de un mismo autor operando en ambas
ciudades. Esta elección introduce además un cambio de dominio entre ciudades,
que deberá tenerse en cuenta al interpretar los resultados de la evaluación por
autor (\ref{sec:results}). La asimetría entre entrenamiento y evaluación es,
por tanto, intrínseca al problema de aprendizaje de métricas.

Las métricas ARI, NMI y F1 por pares se aplican directamente sobre los
\textit{clusters} obtenidos y las etiquetas de estilo y autoría, sin resolver
explícitamente el problema de asignación clúster-clase: Las tres son
invariantes ante permutaciones de etiquetas y comparan las relaciones de
coagrupamiento entre pares de puntos. Para los algoritmos que pueden etiquetar
puntos como ruido ($-1$) se reportan dos lecturas de cada métrica. En la
lectura por defecto los puntos de ruido se tratan como un único clúster
adicional, de modo que influyen en la métrica en lugar de ignorarse
silenciosamente. En la lectura complementaria (\texttt{*\_no\_noise}) esas
instancias se descartan antes de puntuar, de forma que los algoritmos de
densidad y los particionales se comparan sobre el subconjunto que cada uno
realmente agrupa. Esta doble lectura es asimétrica respecto al tratamiento del
ruido en las métricas intrínsecas, donde el ruido siempre se excluye y se
complementan con valores como el porcentaje de ruido.

Para evaluar la calidad de la búsqueda por similitud se emplean también las
etiquetas disponibles según la evaluación: las métricas P@$k$ (precisión en las
$k$ primeras posiciones), MAP@$k$ (precisión media promedio), MRR (rango
recíproco medio) y R@$k$ (\textit{recall} a $k$) se calculan por
consulta y se promedian sobre todas las consultas realizadas para obtener una
media global de la calidad de la búsqueda por similitud. Como R@$k$ acota su
denominador al máximo de vecinos relevantes recuperables en $k$ posiciones
($\min(k,\,|\text{clase}|-1)$), cuando las clases son grandes frente a $k$
(como en la evaluación de estilo, con cuatro clases densas) ese denominador se
satura en $k$ y R@$k$ coincide numéricamente con P@$k$. En ese caso se omite de
las tablas por redundancia y solo se incluye cuando ambas divergen.

\subsubsection{Salida y agregación de resultados}
\label{subsubsec:salida}
Cada lote de \textit{benchmarks} produce un único fichero JSON con dos
secciones: un bloque con información sobre el entorno de ejecución (versión de
Python, librerías principales, \textit{hardware} utilizado), y otro bloque con
una entrada por configuración que contiene todas las métricas y tiempos
capturados. Para las etapas repetidas se incluyen tanto los valores medios como
las desviaciones típicas y el número de muestras agregadas en cada caso, de
manera que el análisis posterior pueda distinguir diferencias dentro del ruido
de medición.

\subsection{Componentes específicos para grafiti}
\label{subsec:entrenamiento}
Más allá del catálogo genérico, el sistema incorpora dos componentes ajustados
al dominio del grafiti: un detector que segmenta los grafitis dentro de cada
fotografía y unas cabezas de proyección que adaptan las representaciones. Ambos
se entrenan sobre los subconjuntos anotados descritos en \ref{subsec:dataset},
disjuntos del banco de evaluación, lo que evita sesgar las métricas por
solapamiento entre entrenamiento y evaluación.

\subsubsection{Detector de grafiti}
\label{subsubsec:segmentador}
\label{subsubsec:entrenamiento-detector}
La etapa de segmentación localiza los grafitis para que el extractor opere sobre
los recortes y no sobre la escena completa. Se implementan dos variantes: una
sin segmentación, que entrega la fotografía completa, y otra basada en un modelo
de YOLO~\cite{ultralytics_ultralytics_2024} (oficial o ajustado sobre grafiti)
que detecta las cajas de cada grafiti. Las detecciones se fusionan por
solapamiento y por un criterio de proximidad que une cajas adyacentes
(consolidando grafitis contiguos en una misma pared), con una salvaguarda que
rechaza uniones que superen una fracción del área de la imagen, truncado
opcional a las $K$ cajas de mayor confianza y un margen configurable alrededor
de cada recorte. Cuando no hay ninguna detección, la imagen se emite completa
como caja única o se descarta. La calidad del detector se mide con las métricas
de detección estándar (precisión, \textit{recall} y mAP a distintos umbrales de
intersección sobre unión, IoU).

El detector se entrena sobre las $447$ fotografías anotadas de ambos bancos
(unas $650$ cajas), partiendo de los pesos preentrenados sobre COCO y ajustando
una única clase, \enquote{grafiti}. Dado el reducido número de ejemplos se
aplica una aumentación agresiva. Se conservan los hiperparámetros por defecto de
la implementación de referencia salvo el número de épocas, el tamaño de
\textit{batch} y la resolución, fijados para encajar en la memoria de la GPU
(\ref{subsec:entorno}), con parada temprana sobre la mAP@$0.5$ de validación.

\subsubsection{Cabezas de proyección ajustadas al dominio}
\label{subsubsec:cabezas}
\label{subsubsec:entrenamiento-cabezas}
Los modelos de fundación generalistas producen representaciones transferibles
pero no óptimas para el grafiti. Para acortar esa distancia se sigue un esquema
de \textit{transfer learning}~\cite{pan_survey_2010,howard_universal_2018} en el
que el \textit{backbone} permanece congelado y solo se entrena una cabeza ligera
(secuencia Lineal $\rightarrow$ ReLU $\rightarrow$ Lineal) que proyecta los
\textit{embeddings} base a un espacio de $128$ dimensiones (DINOv2 $384\to128$,
MobileNetV3 $1280\to128$), normalizado a norma unitaria ($L_2$) para que sean
directamente comparables mediante distancia coseno
(figura~\ref{fig:cabeza-proyeccion}). Se entrenan cuatro cabezas, de autoría y
de estilo sobre MobileNetV3 y DINOv2: las de autoría acercan en el espacio
latente los grafitis de un mismo autor, mientras que las de estilo organizan el
espacio según las cuatro categorías estilísticas.

\begin{figure}[h!]
  \centering
  \shorthandoff{<>}%
  \begin{tikzpicture}[
    >=Stealth, node distance=6mm and 6mm,
    block/.style={draw, rounded corners, minimum height=9mm, minimum width=14mm, inner sep=3pt, align=center, font=\small},
    frozen/.style={block, fill=gray!20, dashed},
    trained/.style={block, fill=blue!12},
    vec/.style={font=\footnotesize, align=center},
  ]
    \node[font=\small] (img) {imagen};
    \node[frozen, right=of img] (bb) {Backbone\\congelado};
    \node[trained, right=of bb] (l1) {Lineal};
    \node[trained, right=of l1] (relu) {ReLU};
    \node[trained, right=of relu] (l2) {Lineal};
    \node[block, fill=green!12, right=of l2] (norm) {$L_2$};
    \node[vec, right=of norm] (e) {embedding};

    \draw[->] (img)--(bb); \draw[->] (bb)--(l1);
    \draw[->] (l1)--(relu); \draw[->] (relu)--(l2); \draw[->] (l2)--(norm);
    \draw[->] (norm)--(e);

    \node[draw, dotted, fit=(l1)(relu)(l2), inner sep=2mm,
          label=below:{\scriptsize Cabeza entrenada}] {};
  \end{tikzpicture}%
  \shorthandon{<>}%
  \caption[Cabeza de proyección ajustada a grafiti]{Cabeza de proyección ajustada a grafiti. El \textit{backbone} permanece
    congelado y solo se entrena una secuencia ligera de capas para luego normalizar
    la salida a norma unitaria.}
  \label{fig:cabeza-proyeccion}
\end{figure}

Las cuatro cabezas se entrenan sobre la totalidad del subconjunto anotado, sin
partición de validación y con los \textit{backbones} congelados. Las de autoría
emplean la pérdida de tripletes \textit{batch-hard}~\cite{hermans_defense_2017}
(margen $0.3$, Adam con \textit{learning rate} $10^{-4}$, $20$ épocas) y un
muestreador PK de $8$ autores y $4$ recortes por autor, que concentra la pérdida
en los tripletes más informativos de cada \textit{batch}. Las de estilo emplean
la pérdida \textit{Supervised Contrastive}~\cite{khosla_supervised_2021}
($T=0.07$, Adam $10^{-4}$, \textit{scheduler} cosenoidal y $60$ épocas) con un
muestreador balanceado y parada temprana (paciencia de $15$ épocas), conservando
la cabeza que minimiza la pérdida.
